\documentclass[11pt,reqno]{amsart}
\usepackage[T1]{fontenc}\usepackage{lmodern}
\usepackage{amsmath,amssymb,amsthm,mathtools}
\usepackage[margin=1in]{geometry}\usepackage{microtype}\usepackage{booktabs}\usepackage{xcolor}
\usepackage[colorlinks=true,linkcolor=blue!55!black]{hyperref}
\theoremstyle{plain}\newtheorem{proposition}{Proposition}\newtheorem{corollary}[proposition]{Corollary}
\theoremstyle{remark}\newtheorem{remark}[proposition]{Remark}
\newcommand{\CM}{C_M}\newcommand{\CF}{C_F}\newcommand{\bp}{\beta+p}
\newcommand{\phio}{\varphi^{\circ}}\newcommand{\CMbar}{\bar C_M}\newcommand{\eqdef}{\vcentcolon=}
\setlength{\parskip}{3pt}
\begin{document}
\title[Theory behind the calibration]{Theory Behind the Calibration:\\
Break-Even Separability and Threshold Insensitivity}
\author{Team A1 --- Role 1}
\date{Supporting \S7}
\begin{abstract}
The calibration section states two facts it does not prove. First, that the break-even
Bitcoin price is \emph{additively separable}, with machine efficiency $\varepsilon$ setting
a power-bill floor and the hardware premium depending only on $q/T$ and ``not on
$\varepsilon$ at all\dots the two effects cancel exactly.'' Second, that the entry
threshold $\CMbar$ is insensitive to the outage penalty once $\bp\gg c+\CF$. Both are used
to license headline numbers---the \$65k cheapest entry point, the \$43.6k power-bill floor,
and a deliberately conservative $\bp=\$200$/MWh. We prove both, and verify that the proved
formulas reproduce every cell of the calibration table.
\end{abstract}
\maketitle

\section{Setup}
A machine of efficiency $\varepsilon$ (J/TH) earns, per MWh consumed,
$R(\varepsilon)=\chi P_{\mathrm{BTC}}/\varepsilon$ with
$\chi\eqdef b\cdot3.6\times10^{9}/(t_B H)$ fixed by network conditions ($b$ the block
reward, $t_B$ the block interval, $H$ the network hash rate). Hardware at $q$ \$/TH
amortized over $T$ years costs, per MW per hour,
\[
  \CM(\varepsilon,q,T)\;=\;\frac{q\cdot 10^{6}/\varepsilon}{8760\,T},
\]
since one MW packs $10^{6}/\varepsilon$ TH/s. Entry occurs iff $\CM<\CMbar=(R-p)\phio$.

\section{Break-even separability}

\begin{proposition}[The break-even price is additively separable, and the hardware premium
does not depend on efficiency]\label{prop:sep}
Entry occurs iff $P_{\mathrm{BTC}}>P^{*}$, where
\[
  \boxed{\;P^{*}(\varepsilon,q,T)\;=\;\underbrace{\frac{p}{\chi}\,\varepsilon}_{\text{power-bill floor}}
  \;+\;\underbrace{\frac{10^{6}}{8760\,\phio\,\chi}\cdot\frac{q}{T}}_{\text{hardware premium}}\;}
\]
The two terms are additively separable; the hardware premium depends on $q/T$ alone and is
\emph{exactly} independent of $\varepsilon$.
\end{proposition}
\begin{proof}
Entry is $(R-p)\phio>\CM$, i.e.
$\big(\chi P_{\mathrm{BTC}}/\varepsilon-p\big)\phio>q\cdot10^{6}/(\varepsilon\,8760\,T)$.
Multiply both sides by $\varepsilon/\phio>0$:
\[
  \chi P_{\mathrm{BTC}}-p\,\varepsilon\;>\;\frac{q\cdot10^{6}}{8760\,\phio\,T},
\]
and divide by $\chi>0$. The $\varepsilon$ appearing on the right has cancelled because
\emph{both} sides of the entry condition are proportional to $1/\varepsilon$: a more
efficient machine earns proportionally more per MWh ($R\propto1/\varepsilon$) and also
packs proportionally more terahash into a MW, so it costs proportionally more per MW
($\CM\propto1/\varepsilon$). The ratio that decides entry is efficiency-free; only the
\emph{power bill}, which is charged per MWh rather than per terahash, retains $\varepsilon$.
\end{proof}

\begin{corollary}[Efficiency competes down the floor; hardware markets set the premium]
$\partial P^{*}/\partial\varepsilon=p/\chi>0$ and $\partial P^{*}/\partial q>0$,
$\partial P^{*}/\partial T<0$, with $\partial^{2}P^{*}/\partial\varepsilon\,\partial(q/T)=0$.
Efficiency gains and hardware-price declines therefore lower the break-even price through
\emph{independent} channels, and a firm cannot trade one against the other.
\end{corollary}

\section{Numerical validation}
Taking $p=\$50$/MWh and $\phio=145/155$, and calibrating $\chi$ from the reported
efficiency coefficient, we recover $\chi=0.019739$ and
\[
  \frac{p}{\chi}=2533.0,\qquad \frac{10^{6}}{8760\,\phio\,\chi}=6181.9,
\]
against the stated $2{,}533$ and $6{,}183$. Proposition~\ref{prop:sep} then reproduces the
calibration table exactly:

\begin{center}\small
\begin{tabular}{llrrrr}
\toprule
& $q$ (\$/TH) & $T=2$ & $T=3$ & $T=4$ & $T=5$\\
\midrule
Fleet average ($\varepsilon=17.2$) & 25 & \$121k & \$95k & \$82k & \$74k\\
                                   & 35 & \$152k & \$116k & \$98k & \$87k\\
                                   & 40 & \$167k & \$126k & \$105k & \$93k\\
Frontier ($\varepsilon=13.5$)      & 25 & \$111k & \$86k & \$73k & \textbf{\$65k}\\
                                   & 35 & \$142k & \$106k & \$88k & \$77k\\
                                   & 40 & \$158k & \$117k & \$96k & \$84k\\
\bottomrule
\end{tabular}
\end{center}

The $q\to0$ limit gives the power-bill floor $2533\times17.2=\$43.6$k at fleet-average
efficiency, and the cheapest cell---frontier hardware, favourable prices, patient
five-year amortization---is $\$65$k, both as stated.

\section{Threshold insensitivity}

\begin{proposition}[$\CMbar$ is insensitive to the outage penalty once $\bp\gg c+\CF$]
\label{prop:insens}
With $\phio=1-\CF/(\bp-c)$ and $\CMbar=(R-p)\phio$,
\[
  \frac{\partial\,\CMbar}{\partial(\bp)}=\frac{(R-p)\,\CF}{(\bp-c)^{2}}>0,
  \qquad \CMbar\uparrow(R-p)\ \text{ as }\ \bp\to\infty,
\]
so the threshold is increasing but at a rate falling like $(\bp)^{-2}$, and it is bounded
above by the miner margin $R-p$ regardless of how the outage penalty is set.
\end{proposition}
\begin{proof}
Differentiate $\phio$ and multiply by $R-p$; the limit is $\phio\to1$.
\end{proof}

At the calibration ($R-p=19$, $c=45$, $\CF=10$):

\begin{center}\small
\begin{tabular}{rrrr}
\toprule
$\bp$ & $\phio$ & $\CMbar$ & $\partial\CMbar/\partial(\bp)$\\
\midrule
100 & 0.8182 & 15.545 & 0.0628\\
200 & 0.9355 & 17.774 & 0.0079\\
400 & 0.9718 & 18.465 & 0.0015\\
1000 & 0.9895 & 18.801 & 0.0002\\
\bottomrule
\end{tabular}
\end{center}

At $\bp=\$200$/MWh a further \$100/MWh of assumed outage cost moves $\CMbar$ by under
\$0.80/MWh, against a threshold near \$18/MWh. Proposition~\ref{prop:insens} is what
licenses the deliberately conservative choice of $\bp$: the entry margin is a fact about
$R-p$, not about how the regulator values lost load.

\begin{remark}[What this does and does not settle]
These two propositions validate the calibration's \emph{internal} logic---the functional
form of $P^{*}$ and the robustness of $\CMbar$. They say nothing about the input values.
In particular the tariff $p=\$50$/MWh is an assumption, not a measurement; rebuilding it
from utility-level data moves $\CMbar=(R-p)\phio$ one-for-one and is the single most
consequential empirical check on the entry threshold.
\end{remark}

\begin{remark}[Margin: this section is the miner's $R-p$, deliberately]
The break-even here uses the miner margin $R-p$, whereas the capacity theory (Weeks 2--4)
uses the utility margin $R-c$. This is intentional, not a discrepancy: entry is a
\emph{third-party miner's} decision---a miner earns $R$ from computing and pays the tariff
$p$---so $\CMbar=(R-p)\phio$ is the miner's entry margin, while the utility-operated capacity
program earns $R$ and pays generation cost $c$. The two market structures, and the third
($R-c-\tau$ in the deregulated instrument comparison of \S5.3), are reconciled in the
project scope note; \S5.2's decentralization theorem links them.
\end{remark}

\end{document}
